% =========================================================
% GUÍA BÁSICA DE UNIX: LINUX Y macOS
% Compilación: pdfLaTeX
% =========================================================

\documentclass[11pt,a4paper]{article}

% ---------------------------------------------------------
% PAQUETES
% ---------------------------------------------------------
\usepackage[utf8]{inputenc}
\usepackage[T1]{fontenc}
\usepackage[spanish,es-nodecimaldot]{babel}
\usepackage{lmodern}
\usepackage[a4paper,margin=2.5cm]{geometry}
\usepackage{xcolor}
\usepackage{listings}
\usepackage{booktabs}
\usepackage{longtable}
\usepackage{enumitem}
\usepackage{hyperref}
\usepackage{fancyhdr}
\usepackage{tcolorbox}
\usepackage{pmboxdraw}
\usepackage{menukeys}

% ---------------------------------------------------------
% CONFIGURACIÓN DE HYPERREF
% ---------------------------------------------------------
\hypersetup{
    colorlinks=true,
    linkcolor=blue!60!black,
    urlcolor=blue!60!black
}

% ---------------------------------------------------------
% CONFIGURACIÓN DE LISTINGS
% ---------------------------------------------------------
\definecolor{codebg}{RGB}{245,245,245}
\definecolor{codecomment}{RGB}{90,90,90}

% pdfLaTeX + listings necesita declarar explícitamente algunos
% caracteres Unicode que pueden aparecer en ejemplos.
\lstdefinestyle{terminal}{
    backgroundcolor=\color{codebg},
    basicstyle=\ttfamily\small,
    frame=single,
    framerule=0.4pt,
    rulecolor=\color{gray!40},
    columns=fullflexible,
    keepspaces=true,
    showstringspaces=false,
    breaklines=true,
    commentstyle=\color{codecomment},
    language=bash,
    extendedchars=true,
    literate=
        {¿}{{\textquestiondown}}1
        {¡}{{\textexclamdown}}1
        {á}{{\'a}}1
        {é}{{\'e}}1
        {í}{{\'\i}}1
        {ó}{{\'o}}1
        {ú}{{\'u}}1
        {Á}{{\'A}}1
        {É}{{\'E}}1
        {Í}{{\'I}}1
        {Ó}{{\'O}}1
        {Ú}{{\'U}}1
        {ñ}{{\~n}}1
        {Ñ}{{\~N}}1
        {├}{{\pmboxdrawuni{251C}}}1
        {└}{{\pmboxdrawuni{2514}}}1
        {─}{{\pmboxdrawuni{2500}}}1
        {│}{{\pmboxdrawuni{2502}}}1
        {┬}{{\pmboxdrawuni{252C}}}1
        {┤}{{\pmboxdrawuni{2524}}}1
        {┼}{{\pmboxdrawuni{253C}}}1
}
\lstset{style=terminal}

% ---------------------------------------------------------
% CAJAS
% ---------------------------------------------------------
\newtcolorbox{importante}{
    colback=yellow!10,
    colframe=orange!70!black,
    title=\textbf{Importante},
    fonttitle=\bfseries
}

\newtcolorbox{peligro}{
    colback=red!5,
    colframe=red!70!black,
    title=\textbf{¡Cuidado!},
    fonttitle=\bfseries
}

\newtcolorbox{ejercicio}{
    colback=blue!5,
    colframe=blue!60!black,
    title=\textbf{Ejercicio},
    fonttitle=\bfseries
}

% Nota específica para diferencias entre sistemas.
\newtcolorbox{macos}{
    colback=green!5,
    colframe=green!45!black,
    title=\textbf{En macOS},
    fonttitle=\bfseries
}

% ---------------------------------------------------------
% ENCABEZADO Y PIE
% ---------------------------------------------------------
\pagestyle{fancy}
\fancyhf{}
\lhead{Terminal de Unix}
\rhead{Linux y macOS}
\cfoot{\thepage}

% ---------------------------------------------------------
% DOCUMENTO
% ---------------------------------------------------------
\begin{document}


% ---------------------------------------------------------
% PORTADA
% ---------------------------------------------------------
\begin{titlepage}
    \centering


    \vspace*{2cm}


    {\Huge\bfseries Guía básica de la terminal de Linux y macOS (Unix)\par}


    \vspace{1cm}


    {\Large Gestión de directorios y ficheros\par}


    \vspace{1cm}
    
    {\Large Laboratorio de Matemáticas\par}


    \vspace{1cm}


	{Grado en Ingeniería Matemática - Doble Grado en Física e Ingeniería Matemática\par}


    \vspace{1cm}


    {\Large Miguel Beltrá Vidal\par}


    \vfill


    \begin{tcolorbox}[width=0.8\textwidth]
    En esta guía aprenderemos a utilizar los principales comandos
    de la terminal de Linux y macOS para movernos por el sistema de archivos
    y gestionar directorios y ficheros.
    \end{tcolorbox}


    \vfill


    {\large Material práctico para estudiantes\par}


    \vspace{0.5cm}


    {\large \today\par}


\end{titlepage}


\tableofcontents
\newpage


% =========================================================
\section{Introducción}
\noindent Esta guía está pensada para trabajar indistintamente con
\textbf{Linux y macOS}. Ambos sistemas proporcionan una terminal con
comandos de estilo Unix, por lo que la mayoría de las órdenes que veremos
son las mismas en ambos sistemas.

\begin{macos}
En macOS puedes abrir la aplicación \textbf{Terminal} desde
\textit{Aplicaciones $\rightarrow$ Utilidades $\rightarrow$ Terminal}.
También puedes utilizar otra aplicación de terminal, como iTerm2. Los
comandos explicados en esta guía se ejecutan de la misma forma.
\end{macos}


% =========================================================


La \textbf{terminal} es una herramienta que permite comunicarnos con el sistema operativo mediante comandos de texto. En esta guía aprenderemos a:


\begin{itemize}
    \item Saber en qué directorio estamos.
    \item Listar el contenido de un directorio.
    \item Cambiar de directorio.
    \item Crear directorios (carpetas).
    \item Crear ficheros.
    \item Copiar ficheros y directorios.
    \item Mover y renombrar ficheros y directorios.
    \item Borrar ficheros y directorios.
    \item Visualizar la estructura de directorios mediante
          \texttt{tree}.
\end{itemize}


\begin{peligro}
Algunos comandos, especialmente \texttt{rm}, pueden borrar
información de forma permanente. Antes de ejecutar un comando
de borrado, comprueba siempre dónde estás y qué vas a borrar.
\end{peligro}


% =========================================================
\section{¿Dónde estoy?}
% =========================================================


Para saber en qué directorio nos encontramos utilizamos:


\begin{lstlisting}
pwd
\end{lstlisting}


\noindent \texttt{pwd} significa \textit{print working directory}. Por ejemplo:


\begin{lstlisting}
$ pwd
/home/alumno
\end{lstlisting}


\noindent Esto indica que estamos dentro del directorio
\texttt{/home/alumno}.


% =========================================================
\section{Listar el contenido de un directorio}
% =========================================================


Para ver qué hay dentro del directorio actual utilizamos:


\begin{lstlisting}
ls
\end{lstlisting}


\noindent Por ejemplo:


\begin{lstlisting}
$ ls
Documents  Downloads  Music  Images
\end{lstlisting}


\subsection{Algunas opciones interesantes}


Para obtener información más detallada:


\begin{lstlisting}
ls -l
\end{lstlisting}


\noindent Para mostrar también los archivos ocultos:


\begin{lstlisting}
ls -a
\end{lstlisting}


\noindent Podemos combinar ambas opciones:


\begin{lstlisting}
ls -la
\end{lstlisting}


\begin{ejercicio}
\begin{enumerate}
    \item Abre una terminal.
    \item Comprueba en qué directorio estás con \texttt{pwd}.
    \item Lista su contenido con \texttt{ls}.
    \item Utiliza \texttt{ls -l}.
    \item Utiliza \texttt{ls -la}.
\end{enumerate}
\end{ejercicio}


% =========================================================
\section{Cambiar de directorio}
% =========================================================


Para desplazarnos entre directorios utilizamos:


\begin{lstlisting}
cd
\end{lstlisting}


\noindent \texttt{cd} significa \textit{change directory}. Por ejemplo, si tenemos un directorio llamado
\texttt{Documentos}:


\begin{lstlisting}
cd Documentos
\end{lstlisting}


\noindent Podemos comprobar que hemos cambiado de directorio:


\begin{lstlisting}
pwd
\end{lstlisting}


\subsection{Volver al directorio anterior}


Para subir un nivel en la estructura de directorios:


\begin{lstlisting}
cd ..
\end{lstlisting}


\noindent Por ejemplo, si estamos en:


\begin{lstlisting}
/home/alumno/Documentos
\end{lstlisting}


\noindent Después de ejecutar:


\begin{lstlisting}
cd ..
\end{lstlisting}


\noindent estaremos en:


\begin{lstlisting}
/home/alumno
\end{lstlisting}


\subsection{Ir al directorio personal}


Si queremos ir al directorio personal, podemos utilizar:

\begin{lstlisting}
cd ~
\end{lstlisting}


\noindent También podemos utilizar simplemente:


\begin{lstlisting}
cd
\end{lstlisting}


\subsection{Ir al directorio raíz}


El directorio raíz de Linux se representa mediante \texttt{/}.


\begin{lstlisting}
cd /
\end{lstlisting}


\begin{ejercicio}
Partiendo de tu directorio personal:


\begin{enumerate}
    \item Entra en \texttt{Documentos}.
    \item Comprueba dónde estás con \texttt{pwd}.
    \item Vuelve al directorio anterior con \texttt{cd ..}.
    \item Entra de nuevo en \texttt{Documentos}.
    \item Vuelve a tu directorio personal utilizando
          \texttt{cd \textasciitilde}.
\end{enumerate}
\end{ejercicio}


% =========================================================
\section{Crear directorios}
% =========================================================


Para crear una carpeta utilizamos:


\begin{lstlisting}
mkdir nombre_carpeta
\end{lstlisting}


\noindent \texttt{mkdir} significa \textit{make directory}. Por ejemplo, si queremos crear:


\begin{lstlisting}
mkdir practica
\end{lstlisting}


Podemos comprobar que se ha creado:


\begin{lstlisting}
ls
\end{lstlisting}


\subsection{Crear varias carpetas}


También podemos crear varias carpetas en un único comando:


\begin{lstlisting}
mkdir documentos imagenes programas
\end{lstlisting}


\begin{ejercicio}
Dentro de tu directorio personal crea la siguiente estructura:




\begin{lstlisting}
linux/
├── documentos/
├── imagenes/
├── programas/
└── ejercicios/
\end{lstlisting}




Puedes hacerlo mediante:


\begin{lstlisting}
mkdir linux
cd linux
mkdir documentos imagenes programas ejercicios
\end{lstlisting}
\end{ejercicio}


% =========================================================
\section{Crear ficheros}
% =========================================================


Una forma sencilla de crear un fichero vacío con el nombre \texttt{nombre\_fichero} es mediante:


\begin{lstlisting}
touch nombre_fichero
\end{lstlisting}


\noindent Por ejemplo:


\begin{lstlisting}
touch notas.txt
\end{lstlisting}


\noindent Comprobamos que existe:


\begin{lstlisting}
ls
\end{lstlisting}


\noindent Podemos crear varios ficheros:


\begin{lstlisting}
touch uno.txt dos.txt tres.txt
\end{lstlisting}


\begin{ejercicio}
Dentro del directorio \texttt{documentos}, crea estos ficheros:


\begin{lstlisting}
apuntes.txt
tareas.txt
examen.txt
\end{lstlisting}


Comprueba que se han creado utilizando:


\begin{lstlisting}
ls
\end{lstlisting}
\end{ejercicio}


% =========================================================
\section{Copiar ficheros}
% =========================================================


Para copiar un fichero utilizamos:


\begin{lstlisting}
cp origen destino
\end{lstlisting}


\noindent Por ejemplo:


\begin{lstlisting}
cp apuntes.txt copia_apuntes.txt
\end{lstlisting}


\noindent Ahora tendremos:


\begin{lstlisting}
apuntes.txt
copia_apuntes.txt
\end{lstlisting}


\noindent También podemos copiar un fichero a otro directorio:


\begin{lstlisting}
cp apuntes.txt ../
\end{lstlisting}


\noindent En este caso, \texttt{../} representa el directorio padre.


\subsection{Copiar un directorio}


Para copiar un directorio completo necesitamos utilizar la opción \texttt{-r}:


\begin{lstlisting}
cp -r documentos documentos_copia
\end{lstlisting}


\noindent La opción \texttt{-r} significa que se copiará el directorio de forma \textbf{recursiva}, incluyendo su contenido.


% =========================================================
\section{Mover y renombrar}
% =========================================================


Para mover un fichero utilizamos:


\begin{lstlisting}
mv origen destino
\end{lstlisting}


\noindent Por ejemplo:


\begin{lstlisting}
mv apuntes.txt ../
\end{lstlisting}


\noindent El fichero \texttt{apuntes.txt} dejará de estar en el directorio actual y aparecerá en el directorio padre.


\subsection{Mover un fichero a otra carpeta}


Por ejemplo:


\begin{lstlisting}
mv examen.txt ../examenes/
\end{lstlisting}


\subsection{Renombrar un fichero}


El comando \texttt{mv} también sirve para cambiar el nombre:


\begin{lstlisting}
mv tareas.txt tareas_linux.txt
\end{lstlisting}


\noindent No estamos creando un fichero nuevo: simplemente estamos cambiando su nombre.


\subsection{Renombrar un directorio}


También podemos hacer:


\begin{lstlisting}
mv imagenes fotos
\end{lstlisting}


% =========================================================
\section{Borrar ficheros}
% =========================================================


Para borrar un fichero utilizamos:


\begin{lstlisting}
rm nombre_fichero
\end{lstlisting}


\noindent Por ejemplo:


\begin{lstlisting}
rm tres.txt
\end{lstlisting}


\noindent El fichero se eliminará.


\begin{peligro}
\texttt{rm} normalmente no envía el fichero a una papelera. La eliminación puede ser permanente.
\end{peligro}


\noindent Podemos pedir confirmación antes de borrar:


\begin{lstlisting}
rm -i tres.txt
\end{lstlisting}


% =========================================================
\section{Borrar directorios}
% =========================================================


Para borrar un directorio vacío podemos utilizar:


\begin{lstlisting}
rmdir nombre_directorio
\end{lstlisting}


\noindent Por ejemplo:


\begin{lstlisting}
rmdir carpeta_vacia
\end{lstlisting}


\noindent \texttt{rmdir} solamente puede eliminar directorios vacíos.


\subsection{Borrar un directorio y todo su contenido}


Para eliminar un directorio que contiene ficheros u otros directorios:


\begin{lstlisting}
rm -r nombre_directorio
\end{lstlisting}


\noindent Por ejemplo:


\begin{lstlisting}
rm -r ejercicios
\end{lstlisting}


\noindent Esto eliminará el directorio \texttt{ejercicios} y todo lo que
contenga.


\begin{peligro}
\texttt{rm -r} puede eliminar muchos archivos de una sola vez. Antes de utilizarlo, comprueba siempre el contenido y asegúrate de que estás en el directorio correcto.
\end{peligro}


% =========================================================
\section{Visualizar la estructura con \texttt{tree}}
% =========================================================


El comando \texttt{tree} permite visualizar los directorios y ficheros formando una estructura de árbol. Por ejemplo, tras escribir:


\begin{lstlisting}
tree
\end{lstlisting}


\noindent Podríamos obtener:


\begin{lstlisting}
.
├── documentos
│   ├── apuntes.txt
│   ├── examen.txt
│   └── tareas.txt
├── ejercicios
├── imagenes
└── programas
\end{lstlisting}


\noindent Aquí, el símbolo \texttt{.} representa el directorio actual.


\subsection{Mostrar solamente directorios}


Podemos utilizar:


\begin{lstlisting}
tree -d
\end{lstlisting}


\noindent Por ejemplo:


\begin{lstlisting}
.
├── documentos
├── ejercicios
├── imagenes
└── programas
\end{lstlisting}


\begin{importante}
El comando \texttt{tree} no forma parte de la instalación mínima de todos los sistemas. En Linux puede instalarse mediante el gestor de paquetes de la distribución; por ejemplo, en Ubuntu se puede utilizar \texttt{sudo apt install tree}. En macOS, una opción habitual es instalarlo mediante Homebrew con \texttt{brew install tree}.
\end{importante}








\section{Limpiar la terminal}


Si queremos limpiar la terminal, podemos utilizar el comando:


\begin{lstlisting}
clear
\end{lstlisting}








\section{Autocompletado}


La mejor forma de escribir el nombre de un directorio o fichero es dejando que sea la propia terminal la que lo escriba. La terminal tiene la función de autocompletado con la tecla \texttt{tab} \keys{\tab}. Por ejemplo, supongamos que nuestro directorio actual tiene la forma:


\begin{lstlisting}
.
├── documentos
├── ejercicios
├── imagenes
└── programas
\end{lstlisting}


Si escribimos 


\begin{lstlisting}
cd do
\end{lstlisting}


\noindent y presionamos la tecla \texttt{tab} \keys{\tab}, entonces se autocompletará como


\begin{lstlisting}
cd documentos
\end{lstlisting}






















\newpage
\appendix
% =========================================================
\section{Resumen de comandos}
% =========================================================


\begin{longtable}{p{0.16\textwidth}p{0.37\textwidth}p{0.35\textwidth}}
\toprule
\textbf{Comando} & \textbf{Función} & \textbf{Ejemplo} \\
\midrule
\endhead


\texttt{pwd} &
Muestra el directorio actual &
\texttt{pwd} \\


\texttt{ls} &
Lista el contenido &
\texttt{ls} \\


\texttt{ls -l} &
Lista con información detallada &
\texttt{ls -l} \\


\texttt{ls -a} &
Incluye archivos ocultos &
\texttt{ls -a} \\


\texttt{cd} &
Cambia de directorio &
\texttt{cd documentos} \\


\texttt{cd ..} &
Sube al directorio padre &
\texttt{cd ..} \\


\texttt{cd \textasciitilde} &
Va al directorio personal &
\texttt{cd \textasciitilde} \\


\texttt{mkdir} &
Crea un directorio &
\texttt{mkdir fotos} \\


\texttt{touch} &
Crea un fichero vacío &
\texttt{touch notas.txt} \\


\texttt{cp} &
Copia un fichero &
\texttt{cp a.txt b.txt} \\


\texttt{cp -r} &
Copia un directorio &
\texttt{cp -r fotos fotos2} \\


\texttt{mv} &
Mueve un fichero/directorio &
\texttt{mv a.txt documentos/} \\


\texttt{mv} &
Renombra &
\texttt{mv a.txt notas.txt} \\


\texttt{rm} &
Borra un fichero &
\texttt{rm notas.txt} \\


\texttt{rmdir} &
Borra un directorio vacío &
\texttt{rmdir temporal} \\


\texttt{rm -r} &
Borra un directorio y su contenido &
\texttt{rm -r temporal} \\


\texttt{tree} &
Muestra la estructura en árbol &
\texttt{tree} \\


\texttt{tree -d} &
Muestra solamente directorios &
\texttt{tree -d} \\


\bottomrule
\end{longtable}




\end{document}